\documentclass[11pt,a4paper]{article}

\usepackage[margin=2.2cm]{geometry}
\usepackage{amsmath,amssymb,bm}
\usepackage{graphicx}
\usepackage{booktabs}
\usepackage{longtable}
\usepackage{siunitx}
\usepackage{hyperref}
\usepackage{float}
\usepackage{subcaption}
\usepackage{xcolor}
\usepackage{enumitem}
\usepackage{microtype}
\usepackage{placeins}

\hypersetup{
    colorlinks=true,
    linkcolor=blue,
    citecolor=blue,
    urlcolor=blue
}

\newcommand{\dchi}{\Delta\chi^2}
\newcommand{\piE}{\boldsymbol{\pi}_{\rm E}}
\newcommand{\pen}{\pi_{{\rm E},N}}
\newcommand{\pee}{\pi_{{\rm E},E}}
\newcommand{\te}{t_{\rm E}}
\newcommand{\uzero}{u_0}
\newcommand{\rhoFS}{\rho}
\newcommand{\Hnull}{H_0}
\newcommand{\Halt}{H_1}

\title{Empirical Calibration and Characterization of the Annual-Parallax LRT\\
for LSST Microlensing Events}
\author{Roman--Rubin hidden-parallax analysis}
\date{\today}

\begin{document}

\maketitle

\begin{abstract}
We study the detectability of annual microlensing parallax in LSST-only finite-source point-lens
(FSPL) light curves using a likelihood-ratio statistic based on the difference in best-fit
$\chi^2$ between a no-parallax model and a parallax model. The analysis uses two dedicated
Monte Carlo populations: an $\Halt$ population generated with annual parallax and an $\Hnull$
population generated without parallax. The $\Hnull$ sample is used to calibrate the null
distribution empirically, while the $\Halt$ sample is used to measure the conditional detection
efficiency and to characterize the events for which parallax is or is not distinguishable.
This document records the statistical hypotheses, simulation and selection pipeline, fitting
policy, calibration tests, diagnostics, and the figure-by-figure interpretation of the current
analysis products.
\end{abstract}

\tableofcontents
\clearpage

% ============================================================
\section{Scientific question}
% ============================================================

The goal is to determine when annual parallax is distinguishable in an LSST microlensing
light curve after the event has already passed the approved pre-fit detectability selection.
The central statistical question is therefore not whether an event is detectable as
microlensing, but whether the data favor a finite-source model with annual parallax over the
corresponding finite-source model without annual parallax.

The primary quantity is
\begin{equation}
    T \equiv \dchi
    =
    \chi^2_{\Hnull}
    -
    \chi^2_{\Halt},
\end{equation}
where both $\chi^2$ values are the final objective values returned by the frozen production
fitting policy.

For perfectly optimized nested models one expects $T\ge 0$. Negative values are therefore not
interpreted as physical evidence against the larger model; they are treated as numerical
nesting violations and audited separately.

% ============================================================
\section{Hypotheses and microlensing models}
% ============================================================

\subsection{Null hypothesis}

The null model is an FSPL light curve without annual parallax:
\begin{equation}
    \Hnull:
    \qquad
    \bm{\theta}_0
    =
    (t_0,u_0,\te,\rhoFS).
\end{equation}

The generating $\Hnull$ simulations are also produced without parallax. No artificial true
values of $\pen$ or $\pee$ are assigned to the $\Hnull$ population.

\subsection{Alternative hypothesis}

The alternative model is an FSPL light curve with annual parallax:
\begin{equation}
    \Halt:
    \qquad
    \bm{\theta}_1
    =
    (t_0,u_0,\te,\rhoFS,\pen,\pee).
\end{equation}

Thus the alternative contains two additional nonlinear parameters,
\begin{equation}
    \Delta k = 2.
\end{equation}

The parallax amplitude is
\begin{equation}
    \pi_{\rm E}
    =
    \left|\piE\right|
    =
    \sqrt{\pen^2+\pee^2}.
\end{equation}

% ============================================================
\section{Likelihood-ratio test}
% ============================================================

For Gaussian photometric uncertainties, the difference in best-fit $\chi^2$ plays the role of
twice the log-likelihood ratio:
\begin{equation}
    T
    =
    \chi^2_{\Hnull}
    -
    \chi^2_{\Halt}.
\end{equation}

Under regularity conditions, Wilks' theorem would suggest
\begin{equation}
    T \sim \chi^2_{\nu=2}
    \qquad
    \text{under }\Hnull.
\end{equation}

For two degrees of freedom,
\begin{equation}
    P(\chi^2_2 > T)
    =
    \exp(-T/2),
\end{equation}
and the asymptotic critical value at false-positive level $\alpha$ is
\begin{equation}
    c_{\alpha,\mathrm{Wilks}}
    =
    -2\ln\alpha.
\end{equation}

However, this problem contains nonlinear parameters, finite bounds, survey sampling,
pre-selection, and numerical optimization. We therefore do not assume that Wilks' theorem is
exact. The final calibration is empirical.

% ============================================================
\section{Monte Carlo populations}
% ============================================================

\subsection{$\Halt$-generated population}

The production run covers the LSSTMONTS catalog domain and applies the approved pre-fit
detectability selection before the LRT fits. The compact production export contains
\begin{equation}
    N_{\Halt}=333\,653
\end{equation}
successfully analyzed $\Halt$-generated events.

The truth model is FSPL with annual parallax. The final compact table stores the LRT statistic,
final $\Hnull$ and $\Halt$ parameters, physical-coordinate covariance information, optimizer
diagnostics, sampling diagnostics, and bookkeeping quantities.

\subsection{$\Hnull$-generated calibration population}

The calibration campaign processed candidate ranks $[0,300000)$ in deterministic order.
The candidate audit is
\begin{align}
    N_{\rm candidates} &= 300\,000,\\
    N_{\rm materialized} &= 103\,937,\\
    N_{\rm not\ detectable} &= 195\,474,\\
    N_{\rm invalid\ catalog\ row} &= 589.
\end{align}

All $103\,937$ materialized events are retained for calibration; there is no truncation to an
exact target of $100\,000$.

% ============================================================
\section{Pre-fit detectability selection}
% ============================================================

The production selection is applied before the LRT fits and defines the population to which
the reported power applies. The frozen configuration requires:
\begin{itemize}[leftmargin=2em]
    \item Rubin photometry only for the pre-fit gate;
    \item at least 10 total photometric points;
    \item at least 3 observed bands;
    \item a peak window of $\pm 1\te$;
    \item at least 5 points in the peak window;
    \item at least one point on each side of the peak within $\pm 1\te$;
    \item at least 6 points above the adopted $3\sigma$ criterion;
    \item minimum Asimov-like $\Delta\chi^2$ per point equal to 2;
    \item the no-parallax model as the reference model for the pre-fit detectability gate.
\end{itemize}

Consequently, every reported $\Halt$ detection fraction is a \emph{conditional} detection
efficiency among events that already pass this pre-fit selection. It is not the detection
fraction over all physical microlensing events in the parent Galactic population.

% ============================================================
\section{Frozen production fitting policy}
% ============================================================

The production policy is intentionally minimal and reproducible.

\subsection{$\Halt$-generated events}

For $\Halt$-generated events, the generating parallax model is fitted once and the nested
no-parallax model is also fitted according to the frozen production policy. The final LRT uses
the final $\chi^2$ values selected by that policy.

\subsection{$\Hnull$-generated events}

For $\Hnull$-generated events:
\begin{enumerate}[leftmargin=2em]
    \item the $\Hnull$ fit is initialized at the true shared parameters
    $(t_0,u_0,\te,\rhoFS)$;
    \item the $\Halt$ fit starts from the exact nested null, with the same shared parameters
    and $(\pen,\pee)=(0,0)$;
    \item all fitted parameters remain free during optimization;
    \item there is one TRF fit for each hypothesis;
    \item no global multistart, differential evolution, continuation, rescue, or projected
    fallback is used in this calibration production.
\end{enumerate}

The covariance matrices stored in the compact exports are expressed in physical parameter
coordinates. The $\Hnull$ covariance order is
$(t_0,u_0,\te,\rhoFS)$ and the $\Halt$ covariance order is
$(t_0,u_0,\te,\rhoFS,\pen,\pee)$.

% ============================================================
\section{Empirical null calibration}
% ============================================================

The empirical critical value is defined by the upper tail of the simulated null population:
\begin{equation}
    \widehat c_\alpha
    =
    \widehat F_0^{-1}(1-\alpha),
\end{equation}
with the implementation choosing an observed null value at or above the requested quantile.

The measured thresholds are:
\begin{table}[H]
\centering
\begin{tabular}{rrrr}
\toprule
$\alpha$ & $c_{\alpha,\chi^2_2}$ & $\widehat c_{\alpha,\Hnull}$ &
$\widehat{\mathrm{FPR}}(c_{\chi^2_2})$\\
\midrule
$5\times10^{-2}$ & 5.991 & 5.742 & 0.043892\\
$10^{-2}$        & 9.210 & 8.946 & 0.008630\\
$10^{-3}$        & 13.816 & 13.160 & 0.000741\\
$10^{-4}$        & 18.421 & 17.867 & 0.000077\\
\bottomrule
\end{tabular}
\caption{Asymptotic and empirically calibrated critical values. The Wilks thresholds are
slightly conservative for the simulated null population.}
\end{table}

The empirical null distribution is therefore close to, but not identical to, $\chi^2_2$.
The production analysis uses empirical calibration rather than relying solely on the
asymptotic approximation.

% ============================================================
\section{Bootstrap uncertainty of the calibrated thresholds}
% ============================================================

Bootstrap resampling of the $\Hnull$ population gives:
\begin{table}[H]
\centering
\begin{tabular}{rrrrr}
\toprule
$\alpha$ & $\widehat c_\alpha$ & bootstrap median & 2.5\% & 97.5\%\\
\midrule
0.05   & 5.742 & 5.741 & 5.696 & 5.790\\
0.01   & 8.946 & 8.935 & 8.845 & 9.033\\
0.001  & 13.160 & 13.160 & 12.719 & 13.540\\
0.0001 & 17.867 & 17.867 & 16.699 & 18.675\\
\bottomrule
\end{tabular}
\caption{Bootstrap uncertainty of empirical critical values. The tail uncertainty grows
rapidly as $\alpha$ decreases.}
\end{table}

At $\alpha=10^{-3}$ the null tail contains about $104$ events, while at
$\alpha=10^{-4}$ it contains only about $10$. We therefore adopt
\begin{equation}
    \boxed{\alpha_{\rm primary}=10^{-3}}
\end{equation}
as the primary working false-positive level for characterization, while retaining the other
levels for sensitivity tests.

% ============================================================
\section{Cross-validated false-positive rate}
% ============================================================

The null sample is split into calibration/test folds. In each split the threshold is estimated
from the calibration subset and the realized false-positive rate is measured on held-out
events.

The mean realized false-positive rates are:
\begin{align}
    \alpha=0.05 &: \quad 0.049973,\\
    \alpha=0.01 &: \quad 0.010006,\\
    \alpha=0.001 &: \quad 0.001020,\\
    \alpha=0.0001 &: \quad 0.000106.
\end{align}

This test is important because it verifies the calibration out of sample rather than simply
checking the same null sample used to define the quantile.

% ============================================================
\section{Conditional stability of the null calibration}
% ============================================================

At the primary threshold $\widehat c_{10^{-3}}=13.159820$, the false-positive rate was examined
as a function of:
\begin{itemize}[leftmargin=2em]
    \item total number of photometric points;
    \item nearest observation to peak in units of $\te$;
    \item number of points within $0.25\te$, $0.5\te$, $1\te$, and $2\te$;
    \item points to the left and right of the peak within $1\te$;
    \item the binary diagnostic \texttt{poor\_peak\_coverage}.
\end{itemize}

Across the sampling bins, the measured false-positive fractions remain of order $10^{-3}$
and their Wilson intervals overlap the target level. The
\texttt{poor\_peak\_coverage=True} subset has an observed false-positive fraction
$7/4408=1.59\times10^{-3}$, but the corresponding confidence interval is broad because only
seven false positives occur in that subgroup. The current result therefore does not establish
a strong sampling-dependent failure of the global threshold.

% ============================================================
\section{Negative LRT values and numerical nesting violations}
% ============================================================

The compact populations contain:
\begin{align}
    N_{\Hnull}(T<0) &= 1722
    \qquad (1.6568\%),\\
    N_{\Halt}(T<0) &= 2978
    \qquad (0.8925\%).
\end{align}

Because the models are nested, negative values cannot represent the ordering of exact global
minima. They are therefore treated as an optimization/nesting diagnostic. Their relation to
peak coverage, optimizer optimality, active bounds, number of evaluations, and success flags
is analyzed separately.

\subsection{Negative events in the null population}

For $\Hnull$-generated events, poor peak coverage is \emph{not} strongly enriched among
negative-LRT cases. The negative fraction is $1.653\%$ for
\texttt{poor\_peak\_coverage=False} and $1.747\%$ for
\texttt{poor\_peak\_coverage=True}. The continuous sampling diagnostics are also very similar
between negative and non-negative events: the median nearest-peak distance is $0.045\te$ versus
$0.055\te$, and the median number of points within $1\te$ is 21 versus 19. Thus sparse peak
sampling is not the dominant explanation for the negative tail under $\Hnull$.

The clearest difference is numerical. The median optimizer optimality rises from 16.35 to
75.60 for the $\Hnull$ fit and from 23.29 to 415.20 for the $\Halt$ fit in negative events.
The $\Hnull$ fit also requires more function evaluations (median 14 versus 9). Conversely, the
$\Halt$ fit uses fewer evaluations in the negative subset (median 10 versus 15). Both optimizer
success flags are nevertheless true for every event. This demonstrates that the binary success
flag alone is too weak a diagnostic: a formally successful local optimization can still leave
the nested ordering numerically inconsistent. Active parameter bounds are not enriched in the
negative subset and therefore do not appear to be the primary cause.

\subsection{Negative events in the parallax-generated population}

The behavior is qualitatively different for $\Halt$-generated events. Poor peak coverage is a
strong risk factor: the negative fraction is $0.800\%$ when the flag is false and $2.989\%$
when it is true, an increase by a factor of approximately 3.7. The negative subset also has
substantially poorer continuous sampling: its median nearest-peak distance is $0.171\te$
compared with $0.055\te$, the median number of points within $0.5\te$ is 5 instead of 10, and
the median number within $1\te$ is 11 instead of 19. The total number of photometric points is
also much lower (median 127.5 versus 311).

However, only 422 of the 2978 negative $\Halt$ events carry the hard
\texttt{poor\_peak\_coverage} flag. The flag therefore identifies an enriched subgroup but does
not explain the majority of negative cases. The continuous sampling variables contain more
information than the binary threshold. The $\Halt$ fit also takes substantially more function
evaluations in the negative subset (median 25 versus 15), while optimizer success remains true
for all events and active bounds are not enriched. The present evidence therefore points to a
combination of weakly constrained/sparsely sampled event geometry and local optimization-basin
issues rather than a simple bound failure or explicit optimizer failure.

These events are not removed from the empirical null distribution merely because they are
negative: the null calibration is intended to characterize the complete frozen analysis
procedure actually used in production. Their existence is nevertheless an important numerical
diagnostic and motivates direct light-curve and trajectory inspection using the saved CHE data
and fitted parameters, without re-simulation or re-fitting.

% ============================================================
\section{False-positive tail under \texorpdfstring{$\Hnull$}{H0}}
% ============================================================

At the empirically calibrated thresholds the null tail contains:
\begin{table}[H]
\centering
\begin{tabular}{rrr}
\toprule
$\alpha$ & $\widehat c_\alpha$ & null events above threshold\\
\midrule
0.05   & 5.742 & 5196\\
0.01   & 8.946 & 1039\\
0.001  & 13.160 & 103\\
0.0001 & 17.867 & 10\\
\bottomrule
\end{tabular}
\caption{False-positive tail of the empirical null distribution.}
\end{table}

The events in this tail are the most important null cases for physical and numerical
inspection, because they are precisely the events that the complete procedure would classify
as parallax detections despite being generated without parallax.

At the primary $\alpha=10^{-3}$ threshold there are 103 such events. They do not show the
signature expected for a simple poor-sampling failure. Their median total number of photometric
points is 324, compared with 310 for the remainder of the null population, and they have
slightly more points within $0.5\te$, $1\te$, and $2\te$. Their median nearest-peak distance is
only mildly larger ($0.059\te$ versus $0.055\te$). The measured false-positive rate for
\texttt{poor\_peak\_coverage=True} is $1.588\times10^{-3}$ compared with
$0.965\times10^{-3}$ for the well-covered subset, but this result is based on only seven false
positives in the poor-coverage group and the Wilson intervals overlap broadly.

The optimizer diagnostics likewise do not reveal a single dominant pathology. The parallax fit
uses more evaluations in the false-positive tail (median 19 versus 15) and has a higher median
optimality value (74.2 versus 24.0), but all fits report optimizer success. Active bounds are
not strongly enriched. These null-tail events should therefore be interpreted primarily as the
expected extreme tail of the complete calibrated procedure unless event-by-event inspection
reveals a reproducible numerical failure mode. In particular, one should not attempt to
``explain away'' all 103 events: a calibration at $\alpha=10^{-3}$ is expected to retain about
one false positive per thousand null events.

% ============================================================
\section{ROC characterization}
% ============================================================

Using $T=\dchi$ as the ranking statistic, the empirical receiver-operating characteristic has
\begin{equation}
    \mathrm{AUC}=0.8983.
\end{equation}

The ROC curve is not used to replace the calibrated false-positive threshold. Its role is to
summarize the global separation of the $\Hnull$ and $\Halt$ populations over the complete
range of possible thresholds.

% ============================================================
\section{LRT versus local Wald significance of parallax}
% ============================================================

For the fitted parallax vector we define the local quadratic statistic
\begin{equation}
    D_{\pi_{\rm E}}^2
    =
    \begin{pmatrix}
    \pen & \pee
    \end{pmatrix}
    C_{\pi_{\rm E}}^{-1}
    \begin{pmatrix}
    \pen\\
    \pee
    \end{pmatrix},
\end{equation}
where $C_{\pi_{\rm E}}$ is the $2\times2$ parallax covariance block from the $\Halt$ fit.

A valid covariance is available for $333\,631$ of $333\,653$ $\Halt$ events. Among events
used in the positive-statistic comparison,
\begin{equation}
    \mathrm{corr}\!\left[
        \log_{10} T,
        \log_{10} D_{\pi_{\rm E}}^2
    \right]
    =
    0.788.
\end{equation}

The strong but imperfect correlation shows that the LRT and the local covariance-based
significance probe related, but non-identical, aspects of the fitted likelihood surface.
Events with strong disagreement between these quantities deserve separate characterization.

% ============================================================
\section{Calibrated detection efficiency under \texorpdfstring{$\Halt$}{H1}}
% ============================================================

Applying the empirical null thresholds to the $\Halt$ population gives:
\begin{table}[H]
\centering
\begin{tabular}{rrrr}
\toprule
$\alpha$ & $\widehat c_\alpha$ & detected $\Halt$ events & conditional power\\
\midrule
0.05   & 5.742  & 258622 & 0.775123\\
0.01   & 8.946  & 243336 & 0.729309\\
0.001  & 13.160 & 229921 & 0.689102\\
0.0001 & 17.867 & 219415 & 0.657614\\
\bottomrule
\end{tabular}
\caption{Conditional parallax-detection efficiency among $\Halt$ events that pass the pre-fit
detectability selection.}
\end{table}

At the primary level,
\begin{equation}
    P_{\Halt}(T>\widehat c_{10^{-3}})
    =
    0.6891.
\end{equation}

This is a population-averaged conditional power, not a universal probability for an arbitrary
microlensing event.

% ============================================================
\section{Event classification for physical characterization}
% ============================================================

At the primary empirical threshold
\begin{equation}
    \widehat c_{10^{-3}} = 13.15982011480088,
\end{equation}
the full $\Halt$ production contains
\begin{align}
    N(T<0) &= 2978,\\
    N(T=0) &= 7,\\
    N(T>0) &= 330\,668.
\end{align}

For the physical characterization below we deliberately restrict to the
\emph{strictly positive} population, $T>0$.  Within that population,
\begin{align}
    0<T<\widehat c_{10^{-3}}
    &: \quad \text{confused with FSPL without parallax},\\
    T\ge \widehat c_{10^{-3}}
    &: \quad \text{parallax detected}.
\end{align}
The corresponding counts are
\begin{align}
    N_{\rm confused} &= 100\,747,\\
    N_{\rm detected} &= 229\,921,
\end{align}
so that the detection fraction conditional on $T>0$ is
\begin{equation}
    P(\mathrm{detected}\mid T>0)=0.695323.
\end{equation}
This number must not be confused with the full $\Halt$ conditional power
$0.689102$, whose denominator also contains the negative and exactly-zero
LRT events.

The compact classified table is
\begin{verbatim}
analysis/lrt/data/h1_lrt_classified_20260921.parquet
\end{verbatim}

% ============================================================
\section{Repository layout for the confirmed characterization}
% ============================================================

All characterization scripts discussed in the following sections live under
\begin{verbatim}
analysis/lrt/characterization/
\end{verbatim}
with numerical outputs under
\begin{verbatim}
analysis/lrt/results/characterization/
\end{verbatim}
and figures under
\begin{verbatim}
analysis/lrt/figures/characterization/
\end{verbatim}

The truth-joined positive-$T$ table used by the later characterization is
\begin{verbatim}
analysis/lrt/results/characterization/18_positive_truth/
    h1_positive_truth_characterization.parquet
\end{verbatim}

The frozen empirical thresholds are stored in
\begin{verbatim}
analysis/lrt/data/h0_empirical_calibration_thresholds_20260921.csv
\end{verbatim}

The following script sequence is the confirmed characterization chain retained
in the repository:
\begin{verbatim}
10_build_h1_parameter_covariance_table.py
11_audit_covariance_geometry.py
12_covariance_geometry_by_lrt_region.py
13_h1_covariance_volume_vs_lrt.py
14_h1_parallax_area_vs_lrt.py
15_h1_offdiagonal_correlations_vs_lrt.py
16_h0_h1_parameter_displacement_vs_lrt.py
17_h1_parameters_uncertainties_vs_lrt.py
18_build_positive_truth_table.py
19_detection_fraction_vs_true_tE.py
20_detection_fraction_true_tE_piE.py
21_piE_amplitude_recovery_vs_truth.py
22_piE_amplitude_uncertainty_mc_vs_truth.py
23_piE_amplitude_mc_coverage_vs_truth.py
24_confused_piE_mc_interval_outcomes.py
25_piE_component_recovery_vs_truth.py
26_piE_component_uncertainty_vs_truth.py
27_piE_component_one_sigma_coverage.py
28_piE_2d_covariance_coverage_vs_truth.py
29_detection_fraction_vs_finite_source_strength.py
\end{verbatim}

Script 30, which introduced a spline/logistic adjustment of the finite-source
trend, is intentionally excluded from the confirmed sequence because its
interpretation was judged insufficiently transparent for the present analysis.

% ============================================================
\section{Covariance geometry and LRT-region characterization}
% ============================================================

\subsection{Physical-coordinate covariance table}

Script
\begin{verbatim}
analysis/lrt/characterization/
    10_build_h1_parameter_covariance_table.py
\end{verbatim}
builds the event-level physical-coordinate covariance characterization.  The
resulting table has $333\,653$ events and 180 columns and is stored at
\begin{verbatim}
analysis/lrt/results/characterization/10_parameter_covariance/
    h1_parameter_covariance_characterization.parquet
\end{verbatim}

\subsection{Dimensionless covariance geometry}

Because raw covariance determinants mix parameters with different units, the
shape analysis uses the correlation matrix $R$ rather than the raw covariance
matrix.  Script
\begin{verbatim}
analysis/lrt/characterization/11_audit_covariance_geometry.py
\end{verbatim}
produces
\begin{verbatim}
analysis/lrt/results/characterization/11_covariance_geometry/
    h1_covariance_geometry.parquet
\end{verbatim}

For the full $\Halt$ population, $329\,974$ correlation matrices are strictly
positive definite.  The median correlation-volume shape factor is
\begin{equation}
    \mathrm{median}\!\left[\sqrt{\det R}\right]
    =4.858\times10^{-4},
\end{equation}
and the median correlation-matrix condition number is approximately
$1.67\times10^4$.  Representative median absolute pair correlations are
\begin{align}
    |\mathrm{corr}(u_0,\te)| &\simeq 0.892,\\
    |\mathrm{corr}(\te,\pen)| &\simeq 0.825,\\
    |\mathrm{corr}(\te,\pee)| &\simeq 0.840,\\
    |\mathrm{corr}(\pen,\pee)| &\simeq 0.770.
\end{align}
These values show that strong parameter degeneracies are a normal feature of
the fitted $\Halt$ population and cannot by themselves be used as a pathology
flag.

\subsection{Negative LRT covariance geometry}

Scripts 12 and 15 show that negative-LRT events are enriched in specific
correlations rather than in uniformly large correlation among every parameter
pair.  Median absolute correlations for negative versus detected events include
\begin{align}
 |\mathrm{corr}(t_0,\pen)| &: 0.737 \ \mathrm{vs}\ 0.462,\\
 |\mathrm{corr}(t_0,\pee)| &: 0.732 \ \mathrm{vs}\ 0.467,\\
 |\mathrm{corr}(t_0,u_0)| &: 0.623 \ \mathrm{vs}\ 0.463,\\
 |\mathrm{corr}(t_0,\te)| &: 0.621 \ \mathrm{vs}\ 0.462,\\
 |\mathrm{corr}(\pen,\pee)| &: 0.867 \ \mathrm{vs}\ 0.748.
\end{align}
In contrast, correlations such as $(u_0,\te)$ and $(\te,\pee)$ are not
enhanced in the negative subset.  At a relative positive-semidefinite tolerance
of $10^{-8}$, about $8.3\%$ of the negative $\Halt$ covariance matrices remain
substantially non-PSD, compared with a much smaller fraction in the positive
population.  This supports an association between negative LRT values and a
more difficult local covariance geometry, without establishing covariance
failure as the sole cause of the negative statistic.

The relevant scripts and result directories are
\begin{verbatim}
analysis/lrt/characterization/
    12_covariance_geometry_by_lrt_region.py
    15_h1_offdiagonal_correlations_vs_lrt.py

analysis/lrt/results/characterization/
    12_covariance_geometry_by_lrt_region/
    15_h1_offdiagonal_correlations_vs_lrt/
\end{verbatim}

\subsection{Correlation-volume and parallax-area scaling with LRT}

Script 13 shows that the dimensionless correlation-volume measure is
non-monotonic with $T$: it increases from the negative-LRT regime to intermediate
positive $T$, then decreases again at very large $T$.  A small
$\det R$ therefore does not uniquely mean ``poor information''; it can also
describe a highly elongated but precisely localized fit.

Script 14 isolates the $2\times2$ parallax covariance block and studies
\begin{equation}
    A_{\rm rel}
    =
    \frac{\sqrt{\det C_{\pi_{\rm E}}}}
         {|\widehat{\boldsymbol{\pi}}_{\rm E}|^2}.
\end{equation}
For positive LRT values this quantity follows an approximately reciprocal
scaling with the LRT statistic.  Representative median
$\log_{10}A_{\rm rel}$ values are approximately $-1.35$ at
$T\simeq95$, $-2.27$ at $T\simeq10^3$, $-3.29$ at $T\simeq10^4$, and
$-4.46$ at $T\simeq1.15\times10^5$.  Thus the parallax error-ellipse area
approximately follows
\begin{equation}
    A_{\rm rel}\propto T^{-1},
\end{equation}
consistent with the local quadratic expectation that a linear fractional
uncertainty scale behaves roughly as $T^{-1/2}$.

The relevant paths are
\begin{verbatim}
analysis/lrt/characterization/
    13_h1_covariance_volume_vs_lrt.py
    14_h1_parallax_area_vs_lrt.py

analysis/lrt/results/characterization/
    13_h1_covariance_volume_vs_lrt/
    14_h1_parallax_area_vs_lrt/
\end{verbatim}

% ============================================================
\section{H0--H1 displacement and fitted-parameter uncertainties}
% ============================================================

Script 16 compares the final $\Hnull$ and $\Halt$ solutions.  Median
displacements for the three LRT regions are
\begin{table}[H]
\centering
\begin{tabular}{lrrrr}
\toprule
region &
$|\Delta t_0|/\te$ &
$|\Delta u_0|$ &
$|\Delta\log_{10}\te|$ &
$|\Delta\log_{10}\rho|$\\
\midrule
$T<0$ & 0.0473 & 0.318 & 0.374 & 1.511\\
$0<T<c_{10^{-3}}$ & 0.0219 & 0.149 & 0.139 & 1.901\\
$T\ge c_{10^{-3}}$ & 0.0895 & 0.232 & 0.221 & 2.504\\
\bottomrule
\end{tabular}
\caption{Median displacement between the final no-parallax and parallax
solutions.}
\end{table}

The detected population can therefore have large $\Hnull$--$\Halt$
morphological displacement because the no-parallax model absorbs true parallax
through the shared parameters.  Negative events are not simply the events with
the largest model-to-model displacement.

Script 17 shows a much clearer distinction in formal uncertainty.  Representative
medians are
\begin{table}[H]
\centering
\begin{tabular}{lrrrr}
\toprule
region &
$\sigma_{u_0}$ &
$\sigma_{\te}/|\widehat{\te}|$ &
$\sigma_\rho/|\widehat\rho|$ &
$\sigma_{\pi_{\rm E},\parallel}/|\widehat{\pi}_{\rm E}|$\\
\midrule
$T<0$ & 5.87 & 8.77 & 89.7 & 9.21\\
$0<T<c_{10^{-3}}$ & 0.975 & 1.27 & 87.2 & 1.64\\
$T\ge c_{10^{-3}}$ & 0.142 & 0.202 & 114.4 & 0.189\\
\bottomrule
\end{tabular}
\caption{Representative median formal uncertainties from the frozen $\Halt$
covariance.  The finite-source parameter $\rho$ is weakly constrained across
all three regions, whereas $u_0$, $\te$, and parallax are especially unstable
for negative-LRT events.}
\end{table}

Repository paths:
\begin{verbatim}
analysis/lrt/characterization/
    16_h0_h1_parameter_displacement_vs_lrt.py
    17_h1_parameters_uncertainties_vs_lrt.py

analysis/lrt/results/characterization/
    16_h0_h1_parameter_displacement_vs_lrt/
    17_h1_parameters_uncertainties_vs_lrt/
\end{verbatim}

% ============================================================
\section{Truth-catalog join and positive-LRT population}
% ============================================================

The production row index was verified to be the zero-based physical row of the
full LSSTMONTS catalog.  The truth join therefore uses \texttt{catalog\_row},
not the catalog field \texttt{Number\_lensing}.  Script
\begin{verbatim}
analysis/lrt/characterization/18_build_positive_truth_table.py
\end{verbatim}
joins all $330\,668$ strict-positive events with the catalog truth and produces
\begin{verbatim}
analysis/lrt/results/characterization/18_positive_truth/
    h1_positive_truth_characterization.parquet
\end{verbatim}

The joined sample has no left-only or right-only failures.  Useful truth
summaries for the positive population are
\begin{align}
    \mathrm{median}(\te) &= 88.20~\mathrm{d},\\
    \mathrm{median}(\pi_{\rm E,true}) &= 0.4779,\\
    \mathrm{median}(D_L) &= 1.215~\mathrm{kpc},\\
    \mathrm{median}(D_S) &= 4.247~\mathrm{kpc},\\
    \mathrm{median}\!\left(\rho/|u_0|\right) &= 1.109\times10^{-3}.
\end{align}
Only 375 positive events satisfy $\rho>|u_0|$, so finite-source strength must
be treated primarily as a continuous variable rather than only through a
source-transit/non-transit split.

As a consistency check, the relative parallax inferred from distances,
$\pi_{\rm rel}$, agrees with $\pi_{\rm E}\theta_{\rm E}$ to a median relative
difference of $5.35\times10^{-8}$.

% ============================================================
\section{Truth-level parallax detectability}
% ============================================================

\subsection{Dependence on event timescale}

Script
\begin{verbatim}
analysis/lrt/characterization/19_detection_fraction_vs_true_tE.py
\end{verbatim}
measures
\begin{equation}
    P(T\ge c_{10^{-3}}\mid T>0,\te^{\rm true}).
\end{equation}
The detection fraction is very small for the shortest events, rises rapidly
between roughly 10 and 100 days, reaches about $0.8$ by $\te\simeq100$ days,
and peaks near $0.86$--$0.87$ around $150$--$200$ days.  At longer timescales
the fraction shows a modest decline/plateau rather than continuing to increase
monotonically.  Thus long $\te$ strongly favors annual-parallax detection but is
not sufficient by itself.

Repository outputs:
\begin{verbatim}
analysis/lrt/results/characterization/19_detection_fraction_vs_true_tE/
analysis/lrt/figures/characterization/19_detection_fraction_vs_true_tE/
\end{verbatim}

\subsection{Joint dependence on timescale and true parallax amplitude}

Script
\begin{verbatim}
analysis/lrt/characterization/20_detection_fraction_true_tE_piE.py
\end{verbatim}
constructs the conditional detection-fraction map in
$(\te^{\rm true},|\pi_{\rm E,true}|)$.  The map shows a strong diagonal
transition: at fixed $\te$, larger true parallax increases detectability, while
at fixed parallax amplitude longer events are generally more detectable.
Short events require substantially larger $\pi_{\rm E}$ to cross the empirical
LRT threshold.  The behavior is qualitatively consistent with the short-event
acceleration scaling
\begin{equation}
    \delta u_{\rm par}\sim
    \pi_{\rm E}\left(\frac{\te}{1~{\rm yr}}\right)^2,
\end{equation}
although cadence, trajectory orientation, peak epoch, and other nuisance
parameters prevent this from being an exact one-parameter boundary.

Repository outputs:
\begin{verbatim}
analysis/lrt/results/characterization/20_detection_fraction_true_tE_piE/
analysis/lrt/figures/characterization/20_detection_fraction_true_tE_piE/
\end{verbatim}

% ============================================================
\section{Parallax-amplitude recovery and Monte Carlo uncertainty}
% ============================================================

\subsection{Amplitude recovery}

Script
\begin{verbatim}
analysis/lrt/characterization/21_piE_amplitude_recovery_vs_truth.py
\end{verbatim}
compares
\begin{equation}
    \Delta\log_{10}\pi_{\rm E}
    =
    \log_{10}
    \left(
    \frac{|\widehat{\boldsymbol{\pi}}_{\rm E}|}
         {|\boldsymbol{\pi}_{\rm E,true}|}
    \right)
\end{equation}
between confused and detected positive-LRT events.

Detected events have a median recovery close to zero over nearly the full
truth range and a comparatively narrow scatter.  Confused events show a strong
positive amplitude bias at low true parallax: near
$\pi_{\rm E,true}\sim0.02$--$0.03$ the median offset is approximately
$0.6$ dex, i.e. a factor of about four in amplitude, and the bias decreases
toward zero as the true parallax increases.

This amplitude bias does not by itself imply a signed bias in either Cartesian
component, because the positive-definite norm of a noisy two-dimensional vector
has a Rice-like upward bias.

Repository outputs:
\begin{verbatim}
analysis/lrt/results/characterization/21_piE_amplitude_recovery_vs_truth/
analysis/lrt/figures/characterization/21_piE_amplitude_recovery_vs_truth/
\end{verbatim}

\subsection{Monte Carlo propagation of the fitted covariance}

For amplitude uncertainties, script
\begin{verbatim}
analysis/lrt/characterization/22_piE_amplitude_uncertainty_mc_vs_truth.py
\end{verbatim}
draws 1024 samples per event from the frozen bivariate normal approximation
\begin{equation}
    (\pen,\pee)^{(k)}
    \sim
    \mathcal N
    \left(
        \widehat{\boldsymbol{\pi}}_{\rm E},
        C_{\pi_{\rm E}}
    \right),
\end{equation}
then transforms every draw to
$|\boldsymbol{\pi}_{\rm E}^{(k)}|$.  This is propagation of the frozen local
covariance, not a new light-curve simulation, refit, or likelihood MCMC.

Of $330\,668$ positive events, $330\,659$ have a valid positive-definite
$2\times2$ parallax covariance for this calculation.  Defining
\begin{equation}
    \sigma_{68,\rm MC}
    =
    \frac{q_{84}-q_{16}}{2},
\end{equation}
the confused population has
\begin{align}
    q_{16}\!\left(\sigma_{68,\rm MC}/\pi_{\rm E,true}\right) &= 0.570,\\
    \mathrm{median}\!\left(\sigma_{68,\rm MC}/\pi_{\rm E,true}\right) &= 1.855,\\
    q_{84}\!\left(\sigma_{68,\rm MC}/\pi_{\rm E,true}\right) &= 13.075,
\end{align}
and only $31.9\%$ have $\sigma_{68,\rm MC}<\pi_{\rm E,true}$.

For detected events,
\begin{align}
    q_{16} &= 0.0530,\\
    \mathrm{median} &= 0.186,\\
    q_{84} &= 0.681,
\end{align}
and $89.5\%$ have $\sigma_{68,\rm MC}<\pi_{\rm E,true}$.

The event-level Monte Carlo quantiles are retained at
\begin{verbatim}
analysis/lrt/results/characterization/
    22_piE_amplitude_uncertainty_mc_vs_truth/
    piE_amplitude_mc_event_intervals.parquet
\end{verbatim}
with the corresponding confirmed figure at
\begin{verbatim}
analysis/lrt/figures/characterization/
    22_piE_amplitude_uncertainty_mc_vs_truth/
    piE_amplitude_uncertainty_mc_vs_truth.png
\end{verbatim}

\subsection{Coverage of the propagated amplitude interval}

Script 23 tests whether the true amplitude lies in the propagated central
interval $[q_{16},q_{84}]$.  Detected events have coverage close to the nominal
68\% over most of the truth range, typically around $0.68$--$0.74$.  Confused
events show severe undercoverage, generally only about $0.15$--$0.4$ over much
of the range.

Script 24 identifies the direction of these failures.  For confused events the
dominant failure mode is
\begin{equation}
    |\boldsymbol{\pi}_{\rm E,true}| < q_{16},
\end{equation}
especially at small true parallax amplitude, where roughly $75$--$85\%$ of
events can fall below the propagated interval.  The opposite failure,
$\pi_{\rm E,true}>q_{84}$, is rare in the same regime.

These results show that a broad interval in the radial amplitude can still
miss the truth systematically because the distribution is centered at a
positive fitted vector and then transformed through a nonlinear norm.

Repository paths:
\begin{verbatim}
analysis/lrt/characterization/
    23_piE_amplitude_mc_coverage_vs_truth.py
    24_confused_piE_mc_interval_outcomes.py

analysis/lrt/figures/characterization/
    23_piE_amplitude_mc_coverage_vs_truth/
    24_confused_piE_mc_interval_outcomes/

analysis/lrt/results/characterization/
    23_piE_amplitude_mc_coverage_vs_truth/
    24_confused_piE_mc_interval_outcomes/
\end{verbatim}

% ============================================================
\section{Truth-vector reconstruction and component-level recovery}
% ============================================================

\subsection{Exact production convention for true parallax components}

The frozen production configuration uses the catalog \texttt{xi} angle with
\begin{verbatim}
angle_basis = galactic_n1n2
component_convention = east_cos_north_sin
\end{verbatim}
The \texttt{alpha} catalog column is retained only as metadata.

The production construction is
\begin{align}
    \pi_{{\rm E},n1} &= \pi_{\rm E}\cos\xi,\\
    \pi_{{\rm E},n2} &= \pi_{\rm E}\sin\xi,
\end{align}
where $n1$ is the local direction of increasing Galactic longitude and $n2$
the local direction of increasing Galactic latitude.  This tangent-plane
vector is then rotated into the local ICRS North/East basis used by pyLIMA.

Script 25 reproduces this exact transformation.  Its maximum relative
amplitude-conservation error over the $330\,668$ positive events is
\begin{equation}
    1.76\times10^{-8},
\end{equation}
validating the reconstructed $(\pen^{\rm true},\pee^{\rm true})$ values.

The event-level truth-component table is
\begin{verbatim}
analysis/lrt/results/characterization/
    25_piE_component_recovery_vs_truth/
    piE_true_components_and_recovery.parquet
\end{verbatim}

\subsection{Recovery of North and East components}

Using the common normalization by total true parallax amplitude,
\begin{equation}
    \epsilon_i
    =
    \frac{
      |\widehat{\pi}_{{\rm E},i}
      -\pi_{{\rm E},i}^{\rm true}|
    }{
      |\boldsymbol{\pi}_{\rm E,true}|
    },
\end{equation}
script 25 gives
\begin{table}[H]
\centering
\begin{tabular}{llrrrr}
\toprule
class & component &
median signed error &
$q_{16}(|\epsilon_i|)$ &
median $|\epsilon_i|$ &
$q_{84}(|\epsilon_i|)$\\
\midrule
confused & $\pen$ & 0.00237 & 0.111 & 0.479 & 1.856\\
confused & $\pee$ & 0.01111 & 0.136 & 0.582 & 2.227\\
detected & $\pen$ & $-7.3\times10^{-5}$ & 0.00958 & 0.0581 & 0.246\\
detected & $\pee$ & $5.37\times10^{-4}$ & 0.0113 & 0.0697 & 0.288\\
\bottomrule
\end{tabular}
\caption{Recovery of the signed parallax components, normalized by the total
true parallax amplitude.}
\end{table}

The signed medians are very close to zero.  Therefore the strong positive bias
seen in $|\widehat{\boldsymbol{\pi}}_{\rm E}|$ for weak/confused events does
not correspond to a comparably strong signed bias in either component.  This
supports the interpretation in terms of noisy Cartesian components followed by
the positive-definite norm.  In this population the East component is, if
anything, slightly less precisely recovered than the North component.

Confirmed figure:
\begin{verbatim}
analysis/lrt/figures/characterization/
    25_piE_component_recovery_vs_truth/
    piE_component_recovery_vs_truth.png
\end{verbatim}

% ============================================================
\section{Calibration of the parallax-component covariance}
% ============================================================

For individual Cartesian components no Monte Carlo transformation is needed:
under the frozen multivariate-normal approximation,
\begin{equation}
    \sigma_N=\sqrt{C_{NN}},
    \qquad
    \sigma_E=\sqrt{C_{EE}}
\end{equation}
are the exact marginal standard deviations.

Script 26 shows that confused events have component uncertainties generally of
order or larger than the total true parallax amplitude, while detected events
typically have
\begin{equation}
    \sigma_{\pi_{{\rm E},i}}/
    |\boldsymbol{\pi}_{\rm E,true}|
    \sim 0.1\text{--}0.25.
\end{equation}
The East uncertainty is modestly larger than the North uncertainty over much
of the population, consistent with the component-recovery result rather than
indicating an unmodeled component-specific failure.

Script 27 evaluates standardized residuals
\begin{equation}
    z_N=
    \frac{\widehat{\pen}-\pen^{\rm true}}{\sqrt{C_{NN}}},
    \qquad
    z_E=
    \frac{\widehat{\pee}-\pee^{\rm true}}{\sqrt{C_{EE}}}.
\end{equation}
For a calibrated Gaussian marginal,
$P(|z_i|\le1)=0.682689$.  The observed one-sigma coverage is generally
approximately $0.75$--$0.88$ over most of the populated parallax range for both
components and both LRT classes.  The marginal covariances are therefore
conservative rather than underestimated.

Script 28 tests the full two-dimensional covariance through
\begin{equation}
    D_{\pi_{\rm E}}^2
    =
    \Delta\piE^T
    C_{\pi_{\rm E}}^{-1}
    \Delta\piE,
\end{equation}
with the nominal 68\% ellipse defined by
\begin{equation}
    D_{\pi_{\rm E}}^2
    \le
    \chi^2_{2,0.68}
    =
    -2\ln(1-0.68)
    \simeq2.279.
\end{equation}
The empirical 2D coverage is also above the nominal 68\% level through most of
the populated truth range, typically around $0.73$--$0.80$.  Confused and
detected events show broadly similar behavior.  Thus the severe radial
undercoverage found for $|\pi_{\rm E}|$ is not evidence that the Cartesian
parallax covariance is too small; it arises primarily from the nonlinear
positive-definite transformation to amplitude, together with conditioning on
the selected LRT populations.

Repository paths:
\begin{verbatim}
analysis/lrt/characterization/
    26_piE_component_uncertainty_vs_truth.py
    27_piE_component_one_sigma_coverage.py
    28_piE_2d_covariance_coverage_vs_truth.py

analysis/lrt/figures/characterization/
    26_piE_component_uncertainty_vs_truth/
    27_piE_component_one_sigma_coverage/
    28_piE_2d_covariance_coverage_vs_truth/

analysis/lrt/results/characterization/
    26_piE_component_uncertainty_vs_truth/
    27_piE_component_one_sigma_coverage/
    28_piE_2d_covariance_coverage_vs_truth/
\end{verbatim}

% ============================================================
\section{Finite-source strength and parallax detectability}
% ============================================================

The first finite-source diagnostic uses the continuous truth quantity
\begin{equation}
    R_{\rm FS}
    =
    \frac{\rho^{\rm true}}{|u_0^{\rm true}|}.
\end{equation}
This is preferable to a binary split because only 375 positive-LRT events
satisfy $R_{\rm FS}>1$.

Script
\begin{verbatim}
analysis/lrt/characterization/
    29_detection_fraction_vs_finite_source_strength.py
\end{verbatim}
measures the raw conditional fraction
\begin{equation}
    P(
      \mathrm{parallax\ detected}
      \mid T>0,R_{\rm FS}
    ).
\end{equation}

For $R_{\rm FS}\lesssim1$ the detection fraction changes only moderately around
the overall positive-LRT value.  It declines from roughly the low/mid-$0.7$
range at very small $R_{\rm FS}$ to a broad minimum around $0.64$--$0.66$ near
$R_{\rm FS}\sim10^{-2}$, then returns toward roughly $0.68$--$0.70$ as
$R_{\rm FS}$ approaches unity.  For source-transit events,
$R_{\rm FS}>1$, the measured detection fraction rises, reaching above $0.8$ in
the largest plotted bins, but the Wilson intervals become broad because this
tail contains few events.

This first result does \emph{not} establish a causal finite-source effect.
Before attributing the source-transit rise to $\rho$ itself, the truth
distributions of $\te$ and $|\pi_{\rm E}|$ must be compared across
$R_{\rm FS}$, because both are already known to be strong determinants of
parallax detectability.

Confirmed outputs:
\begin{verbatim}
analysis/lrt/results/characterization/
    29_detection_fraction_vs_finite_source_strength/

analysis/lrt/figures/characterization/
    29_detection_fraction_vs_finite_source_strength/
    detection_fraction_vs_finite_source_strength.png
\end{verbatim}

% ============================================================
\section{Current characterization conclusions}
% ============================================================

The confirmed characterization supports the following working conclusions:
\begin{enumerate}[leftmargin=2em]
    \item Annual-parallax detectability is controlled strongly by both true
    event timescale and true parallax amplitude, with a diagonal detection
    boundary in the $(\te,\pi_{\rm E})$ plane.

    \item Positive-LRT events below the calibrated threshold are not simply
    detected events with slightly worse precision.  Their fitted parallax
    amplitudes can be strongly biased upward at low true $\pi_{\rm E}$ and
    their propagated radial intervals are very broad.

    \item The strong low-signal bias is primarily a property of the parallax
    \emph{amplitude}.  The signed North and East components have median
    residuals close to zero.

    \item Both the marginal component uncertainties and the full 2D parallax
    covariance ellipse are conservative relative to the realized component
    errors.  The poor radial coverage of $|\pi_{\rm E}|$ should therefore not
    be interpreted as evidence that the Cartesian covariance matrix is too
    small.

    \item Negative-LRT events remain a distinct numerical population.  They
    are associated with poorer sampling and stronger selected degeneracies,
    especially involving $t_0$ and the parallax components, but no single
    covariance or optimizer diagnostic explains all of them.

    \item The first finite-source diagnostic shows only modest variation in
    parallax-detection fraction for $\rho/|u_0|<1$ and a possible rise for
    source-transit events.  The latter must be checked against the associated
    $\te$ and $\pi_{\rm E}$ distributions before assigning a finite-source
    interpretation.
\end{enumerate}

% ============================================================
\section{Interpretation of the current figure set}

The current repository figure set was inspected together with the numerical CSV outputs that
produced it. The automatically generated appendix below contains every current LRT validation
figure, with a figure-specific caption and interpretation. Several general conclusions emerge.

\begin{enumerate}[leftmargin=2em]
    \item The empirical $H_0$ tail is systematically thinner than the $\chi^2_2$ reference over
    the scientifically relevant range. Consequently, Wilks thresholds are slightly conservative.

    \item The bootstrap and held-out tests show that the empirical calibration is stable at
    $\alpha=0.05$, $0.01$, and $10^{-3}$. At $10^{-4}$ the uncertainty is dominated by the fact
    that only about ten null events populate the relevant tail.

    \item No strong monotonic dependence of the null false-positive rate on the tested sampling
    diagnostics is visible. All conditional Wilson intervals are compatible with the global
    target $\alpha=10^{-3}$. The \texttt{poor\_peak\_coverage=True} subset has a larger point
    estimate, but only seven false positives and a correspondingly broad interval.

    \item Negative LRT values have different behavior in the two generated populations. Under
    $H_0$, poor peak coverage produces little change in the negative fraction; the numerical
    diagnostics instead point toward poorer optimizer optimality. Under $H_1$, poor peak coverage
    increases the negative-LRT fraction from about $0.8\%$ to about $3.0\%$, demonstrating a
    strong sampling dependence.

    \item The optimizer-success flags are constant and equal to \texttt{True} in the current
    compact production sample. Figures conditioned only on those flags therefore contain a
    single point and are audit figures rather than scientifically discriminating plots.

    \item The ROC curve shows strong population separation, with AUC $=0.8983$. The low-FPR
    panel is the more relevant operating view because the scientific thresholds lie between
    $10^{-4}$ and $0.05$.

    \item The LRT and covariance-based Wald statistic are strongly correlated but show broad
    scatter, especially at low significance. They should therefore not be treated as
    interchangeable measures of parallax detectability.
\end{enumerate}

For a final paper, the recommended main-text figures are the null survival curve, bootstrap
critical values, cross-validated FPR, a compact conditional-stability summary, the $H_1$ poor
peak-coverage negative-LRT diagnostic, the low-FPR ROC panel, and the LRT--Wald comparison.
The single-category optimizer-success plots are best retained only as reproducibility/audit
material.
\clearpage

% ============================================================
\appendix
\section{Automatically included repository figures}
% ============================================================

\IfFileExists{generated/figure_manifest.tex}{
    \input{generated/figure_manifest.tex}
}{
    \noindent
    \textbf{Figure manifest not found.}
    Run \texttt{python build\_figure\_manifest.py} from this report directory before compiling.
}

\end{document}
